\documentclass[11pt,letterpaper]{article}

\usepackage[T1]{fontenc}
\usepackage[utf8]{inputenc}
\usepackage{lmodern}
\usepackage{microtype}
\usepackage{amsmath,amssymb,mathtools}
\usepackage{booktabs}
\usepackage{array}
\usepackage{enumitem}
\usepackage[top=0.85in,bottom=0.88in,left=0.92in,right=0.92in]{geometry}
\usepackage{xcolor}
\usepackage[
  colorlinks=true,
  urlcolor=blue!55!black,
  linkcolor=blue!45!black,
  citecolor=blue!45!black,
  bookmarksnumbered=true,
  pdfauthor={Refine Technologies, Inc.},
  pdftitle={Report on The Music of the Primes}
]{hyperref}
\usepackage{titlesec}
\usepackage{needspace}

\preto{\subsection}{\needspace{9\baselineskip}}

\setlength{\parindent}{0pt}
\setlength{\parskip}{5pt plus 0.8pt minus 0.4pt}
\setlength{\abovedisplayskip}{6pt plus 1.2pt minus 1pt}
\setlength{\belowdisplayskip}{6pt plus 1.2pt minus 1pt}

\definecolor{rulegray}{gray}{0.72}
\definecolor{locgray}{gray}{0.28}
\definecolor{baraccent}{RGB}{132,38,38}
\definecolor{barbg}{gray}{0.968}

\titleformat{\section}{\large\bfseries}{\thesection.}{0.5em}{}
\titlespacing*{\section}{0pt}{12pt plus 2pt minus 1pt}{3.5pt}

\titleformat{\subsection}{\normalsize\bfseries}{\thesubsection.}{0.45em}{}
\titlespacing*{\subsection}{0pt}{10pt plus 1.5pt minus 1pt}{2pt}

\makeatletter
\newcommand{\msloc}[1]{%
  {\small\color{locgray}\itshape Manuscript reference: #1.}\par\nobreak\vspace{3pt}\@afterindentfalse\@afterheading
}
\newcommand{\issuefinding}[2]{%
  \begingroup
  \setlength{\fboxsep}{0pt}%
  \noindent\colorbox{barbg}{%
    \begin{tabular}{@{}!{\color{baraccent}\vrule width 2pt}@{\hspace{8pt}}p{\dimexpr\linewidth-18pt\relax}@{\hspace{8pt}}}
      \rule{0pt}{11.5pt}\small\textbf{\color{baraccent}#1}\enspace #2\rule[-4.5pt]{0pt}{0pt}
    \end{tabular}%
  }%
  \endgroup
  \par\nobreak\vspace{4.5pt}\@afterindentfalse\@afterheading
}
\makeatother

\newcommand{\Q}{\mathbb{Q}}
\newcommand{\R}{\mathbb{R}}
\newcommand{\Z}{\mathbb{Z}}
\newcommand{\A}{\mathbb{A}}
\newcommand{\C}{\mathbb{C}}
\newcommand{\Zhat}{\widehat{\Z}}
\newcommand{\dd}{\,d}
\DeclareMathOperator{\Res}{Res}
\DeclareMathOperator{\RePt}{Re}
\DeclareMathOperator{\ImPt}{Im}
\DeclareMathOperator{\supp}{supp}

\raggedbottom

\begin{document}
\fontsize{10.5}{12.65}\selectfont

{\LARGE\bfseries Report on \emph{The Music of the Primes}\par}
\vspace{4pt}
{\normalsize Refine Technologies, Inc.\par}
\vspace{2pt}
{\small\color{locgray}\itshape Note: We used the Refine harness and all frontier models to prepare this report.\par}
\vspace{4pt}

\section{Overview, Proof Architecture, and Strategic Assessment}

\textbf{Purpose and scope.} This report evaluates Alessandro Maria Masullo's manuscript \emph{The Music of the Primes} (November~20, 2026), which proposes an unconditional proof of the Generalized Riemann Hypothesis (GRH) for all Hecke $L$-functions of number fields $F$, together with an unconditional bound of $d=[F:\Q]$ on the multiplicities of non-central zeros for finite-order characters (and an analogous bound $m_\chi$ for algebraic characters assuming Schanuel's conjecture; Expectation~1.7.1 and \S14.1.2). The audited PDF is available from \href{https://math216.wordpress.com/wp-content/uploads/2026/11/the-music-of-the-primes.pdf}{Ravi Vakil's Math~216 site} and as an \href{https://raw.githubusercontent.com/Refine-Technologies-Inc/music-of-primes-report-by-refine/main/the-music-of-the-primes.pdf}{archived copy in the report repository}.

\textbf{Executive summary.} A detailed audit identifies ten substantive gaps and errors across the manuscript; in its present form, the manuscript establishes neither GRH nor the non-central zero-multiplicity bound. Two central steps in the final chapters---the global-to-local trace comparison in Section~12.3 and the divisor theorem in Theorem~13.3.1.2---are unwritten, consisting only of headings and editorial placeholders. Eight further foundational or technical claims admit explicit counterexamples, mostly already over $\Q$ for the Riemann zeta function. Most critically, Chapter~11's log-norm truncation collapses every prime shell to zero, removing the geometric mechanism meant to yield strip-width invariance, while Chapters~9--10 leave unproved both a strip approximation theorem and a cluster-based trace-norm bound.

\textbf{Adelic setup and how the proposed proof works.} Because the main text runs to more than five hundred pages before appendices and introduces extensive custom terminology, it helps to state the manuscript's strategy in standard terms of analytic number theory. Let $F$ be a number field of degree $d=[F:\Q]$, $\A_F$ its adele ring, $\A_F^\times$ its idele group, and $|\cdot|\colon \A_F^\times\to\R_{>0}$ the idele norm ($|\xi|=1$ for $\xi\in F^\times$ by the product formula). Write $C_F = \A_F^\times/F^\times$ for the idele class group with Haar measure $\bar\mu$, and $C_F^1 = \{g\in C_F : |g|=1\}$ for its compact norm-one subgroup, normalized to Haar volume $1$. Choosing a continuous norm section $\tau\colon \R_{>0}\to\A_F^\times$ ($|\tau(t)|=t$) decomposes $C_F \cong \R_{>0}\times C_F^1$, parameterizing idele classes as $g = \tau(e^x)k$ with log-norm $x = \log|g|\in\R$ and $k\in C_F^1$. For $F=\Q$, one has $C_\Q^1 \cong \Zhat^\times$; restricting to the \emph{trivial character} of this compact group reduces functions on $C_\Q$ to functions of the single log-norm variable $x\in\R$ and specializes the Hecke $L$-function to $\zeta(s)$.

For a test function $\varphi$ on $C_F$ and a quasi-character $\eta\colon C_F\to\C^\times$, write $M\varphi(\eta) = \int_{C_F}\varphi(g)\eta(g)\dd\bar\mu(g)$ for the global Mellin transform. At each place $v$ of $F$, Weil's explicit formula attaches a local distribution $D_v(\varphi)$ (summing over prime-power norms at finite places and integrating against an archimedean kernel at infinite places); for any finite set of places $S$, set $D_S(\varphi) = \sum_{v\in S} D_v(\varphi)$. With the manuscript's normalization (Theorem~B.4.4.7), Weil's classical explicit formula unconditionally equates the centered sum over nontrivial zeros to
\[
  M\varphi(|\cdot|^{1/2}) + M\varphi(|\cdot|^{-1/2}) - \sum_v D_v(\varphi),
\]
where the two Mellin terms arise from the poles of the Dedekind zeta function at $1$ and $0$, shifted by $1/2$. By itself, Weil's formula leaves the real parts of the zeros unconstrained.

To locate the zeros, the manuscript adapts the spectral-realization paradigm of Weil and Connes. Fix a vertical strip $I$ centered at $\RePt s = 1/2$, and consider a weighted $L^2$ space of functions on $C_F$ whose Mellin transforms are holomorphic on $I$. Dividing out a subspace generated by global adelic zeta transforms---whose Mellin transforms vanish at Hecke zeros---produces a \emph{spectral quotient} intended to encode only the Hecke zeros inside $I$. Convolution by a test function $\varphi$ (chosen to vanish on $C_F^1$ and to avoid prime powers outside a finite set of places $S$) defines a linear operator on this quotient, and the manuscript attempts to compute its trace in two ways:
\begin{enumerate}[leftmargin=1.6em,itemsep=1.5pt,topsep=2pt]
\item \emph{Spectrally (Chapters~9--10):} expanding the operator in zero-evaluation functionals to express the trace (up to sign) as the multiplicity-weighted sum of $M\varphi(c|\cdot|^{-1/2})$ over zeros $c\in I$.
\item \emph{Geometrically (Chapters~11--12):} peeling off individual prime contributions outside $S$ via prime-difference operators (``shells''), lifting the kernel from $C_F$ to $\A_F^\times$, and comparing the residual global trace with a normalized cutoff integral on the $S$-idele class group $J_{F,S}$ to arrive at $D_S(\varphi) - M\varphi(|\cdot|^{1/2}) - M\varphi(|\cdot|^{-1/2})$.
\end{enumerate}
If the same geometric expression matches the spectral trace for arbitrarily narrow strips $I$, zeros off the critical line are excluded. The manuscript attempts to extract GRH from this trace formula by testing against cutoffs of $t^{-z}\mathbf{1}_{t>1}$ (\S1.4), and to deduce the non-central multiplicity bound $\le d$ by passing to a capped trace identity on the critical line (Chapter~13).

\textbf{Strategic classification of the ten issues.} The ten flaws identified in this report carry very different strategic weight. Grouping them into three tiers clarifies where the proposed proof architecture actually breaks down:

\begin{itemize}[leftmargin=1.5em,itemsep=2.5pt,topsep=2pt]
\item \textbf{Fatal obstructions at the heart of the strategy (Section~\ref{sec:fatal}; Issues~1--3).} First, \emph{Chapter~11's prime-shell construction provably loses the primes}: for $\zeta(s)$, cutting off the log-norm variable to any bounded window before taking closures forces every truncated prime-difference space to coincide with the truncated ambient space, so every successive prime shell is identically zero. Second, \emph{Chapter~12's global-to-local trace comparison is unwritten (\S12.3 is a blank heading) and carries full GRH strength}: although the two missing polar terms have a known representation-theoretic source, proving that the \emph{same} geometric expression equals the zero sum inside an arbitrary centered strip already implies GRH via an elementary Laplace-transform argument. Third, \emph{Chapter~13's non-central multiplicity bound rests on an unproved (and unwritten) capped line identity}: positive-width strip traces count critical-line zeros with full multiplicity $m$ rather than $\min(m,d)$, so shrinking strips cannot produce the cap, and Theorem~13.3.1.2 ends immediately after its title.
\item \textbf{Substantial open theorems on the spectral side (Section~\ref{sec:spectral}; Issues~4--5).} In Chapter~9, reversing a mistaken monotone-projection claim leaves unproved a strip-norm approximation theorem showing that every strip function vanishing at the zeros lies in the geometric subspace being divided out. In Chapter~10, clustered simple critical-line zeros of $\zeta(s)$ blow up individual biorthogonal norms and invalidate the coordinatewise trace-norm bound; although cluster cancellation could conceivably preserve nuclearity, Connes's critical-line frequency-grouping proof has no literal open-strip multiplier analogue.
\item \textbf{Locally reparable statements and bypassed deductions (Section~\ref{sec:reparable}; Issues~6--10).} Five further errors falsify printed theorems as stated: an overlap-gluing mismatch in the global converse theorem (Chapter~6); an inverted Fourier cutoff ($X/p$ vs.\ $pX$) in a local Paley--Wiener quotient (Chapter~7); a conflation of continuous ambient support with distributional restriction on coordinate hyperplanes (Chapter~8); an identically zero averaging weight in the idelic kernel lift (Chapter~11); and a divergent archimedean integral in the introductory cutoff deduction (\S1.4). Each admits a plausible local fix or strategic bypass, though downstream properties would still have to be verified.
\end{itemize}

\section{Fatal Obstructions to the Core Proof Strategy}
\label{sec:fatal}

We begin with the three failures that strike directly at the mechanism proposed to prove GRH and bound zero multiplicities. Everything needed to follow them can be stated directly in the log-norm variable on $C_F$.

\subsection{Issue 1: Chapter 11's truncated prime-shell filtration provably loses the primes}
\msloc{Theorem 11.1.2.1(2), p.~480; Proposition 11.2.0.1(2), p.~484; \S11.2.2; used on pp.~495--496 and 519--520}
\issuefinding{False claims.}{Proposition~11.2.0.1(2) falsely asserts that the prime-difference operator has closed range on the weighted global $L^2$ space, and Theorem~11.1.2.1(2) falsely asserts strict inclusion for the closed log-norm truncations along some subsequence; already for $\zeta(s)$, every truncated prime shell collapses to zero.}

To compute the geometric side of the trace formula prime by prime, Chapter~11 filters its global function space by prime-difference operators so that successive quotients (``shells'') isolate individual primes. Work over $F=\Q$ on the trivial character of $C_\Q^1\cong\Zhat^\times$, where the manuscript identifies the completed global space in $x=\log|g|\in\R$ with the polynomially weighted Hilbert space
\[
  H = L^2\bigl(\R, (1+x^2)^{s/2}\dd x\bigr), \qquad s\ge 2.
\]
Up to translation sign, the operator attached to a prime $p$ is $D = I - T$, where $(Th)(x) = h(x - \log p)$ (corresponding on the Mellin side to multiplication by $1 - p^{\pm z}$). Proposition~11.2.0.1(2) first claims that $D(H)$ is closed in $H$, so that \S11.2.2 can form a Hausdorff quotient $H/D(H)$. That already fails on the whole line. Because $T$ is bounded on $H$ and any fixed vector $Th=h$ is $(\log p)$-periodic and therefore $0$ in $H$, the operator $D$ is bounded and injective; if $D(H)$ were closed, the Bounded Inverse Theorem would force $D$ to be bounded below. Yet for any nonzero bump $\psi\in C_c^\infty((1,2))$ and $g_R(x) = \psi(x/R)$ ($R\ge 1$), the supports of $g_R$ and $Tg_R$ lie in $[R, 2R+\log p]$, where $(1+x^2)^{s/2}\asymp R^s$ and the mean value theorem gives $|Dg_R(x)|\le (\log p)\|\psi'\|_\infty/R$. Integrating over this interval of length $O(R)$ yields
\[
  \|g_R\|_H^2 \asymp R^{s+1}, \qquad \|Dg_R\|_H^2 = O(R^{s-1}),
\]
so $\|Dg_R\|_H / \|g_R\|_H = O(R^{-1}) \to 0$ as $R\to\infty$. Thus $D(H)$ is not closed and $H/D(H)$ is non-Hausdorff.

Much worse is what happens when Chapter~11 truncates in log-norm and takes closures to build its prime shells (Theorem~11.1.2.1(2)). Let $V\subset H$ be the manuscript's space of adelic theta test functions on this component; $V$ is invariant under $T$, and imposing the condition at $p$ replaces $V$ by $(I-T)V$. For a bounded interval $J\subset\R$, write $M_J$ for multiplication by $\mathbf{1}_J$. To obtain even the first prime shell, Theorem~11.1.2.1(2) requires $\overline{M_J(I-T)V}$ to be strictly smaller than $\overline{M_J V}$ along some expanding sequence of intervals $J$.

In fact, the two closed spaces are \emph{equal for every bounded interval $J$}. Given any $f\in V$, translation invariance puts $g_N = \sum_{j=0}^{N-1} T^j f$ in $V$ for every $N\ge 1$. Telescoping gives
\[
  M_J(I - T)g_N = M_J f - M_J T^N f \longrightarrow M_J f \qquad \text{in } H \text{ as } N\to\infty,
\]
because the weight $(1+x^2)^{s/2}$ is bounded on $J$ while the unweighted $L^2$ norm of $f$ on the shifted interval $J - N\log p$ tends to zero as that interval escapes to infinity. Combined with the trivial inclusion $(I-T)V\subset V$, this forces $\overline{M_J(I-T)V} = \overline{M_J V}$. Because prime-difference operators commute with translations, the same telescoping argument applies inductively after any finite set of prime conditions has already been imposed: \emph{every successive quotient of these closed, truncated spaces is zero}, and no choice of subsequence can produce strict inclusion.

\textbf{Why truncation loses the primes, and prospect of repair.} Before truncation, periodization $f\mapsto \sum_{k\in\Z} f(u+k\log p)$ annihilates $(I-T)f$ and maps $H$ continuously and surjectively onto $L^2(\R/\Z\log p)$ (indeed for any $s>1$, by Cauchy--Schwarz and zero-extension from a single period). The whole-line quotient $H/\overline{(I-T)H}$ thus genuinely retains periodic data, whereas cutting off to a bounded window $J$ first allows $g_N$ to push the compensating tail $T^N f$ out of view. A revised approach would have to form the quotient \emph{before} localizing, carry a trace through localization, and prove a new comparison theorem relating that localized quotient trace to the global kernel. Note, however, that shell collapse alone neither proves nor disproves the narrower large-prime vanishing claim invoked in Chapter~12.

\subsection{Issue 2: The unwritten Chapter 12 comparison and why strip-invariance forces GRH}
\msloc{Section 12.3, p.~514; Theorem 12.4.1.1, p.~514; Theorem 12.5.1.1, pp.~519--520}
\issuefinding{Invalid inference; unwritten core step.}{In the proof of Theorem~12.5.1.1, the claimed evaluation of the residual global kernel $k_S$ as $\int_{C_F} k_S(x,x)\dd\bar\mu(x) = D_S(\varphi) - M\varphi(|\cdot|^{1/2}) - M\varphi(|\cdot|^{-1/2})$ does not follow from Theorem~12.4.1.1 (which yields $D_S(\varphi)$ alone) or from (11.3.1.4). Section~12.3, assigned to compare the global and local kernels, is unwritten, leaving Theorem~12.5.1.1---a statement equivalent to GRH---unproved.}

Even granting all earlier chapters, consider Chapter~12 with a global test function $\varphi$ on $C_F$ vanishing on $C_F^1$ (so $\varphi(1)=0$) and a finite set of places $S$ chosen so that $\varphi$ avoids prime powers outside $S$. On the local side, Theorem~12.4.1.1 evaluates the normalized diagonal integral of a Paley--Wiener cutoff kernel $k^{\mathrm{PW}}_{S,\varphi,\lambda,\lambda'}$ on the $S$-idele class group $J_{F,S}$. Because its logarithmic divergence is proportional to $\varphi(1)=0$, sending the smoothing width $\varepsilon\to 0$ with $\varepsilon\lambda,\varepsilon\lambda'\to 0$ leaves only
\[
  \lim_{\lambda,\lambda'\to\infty}
  \frac{1}{\bar\mu_{S,1}(J_{F,S}^1)}
  \int_{J_{F,S}} k^{\mathrm{PW}}_{S,\varphi,\lambda,\lambda'}(x,x)\dd\bar\mu_S(x)
  = D_S(\varphi).
\]
In the proof of the Main Trace Formula (Theorem~12.5.1.1, pp.~519--520), the manuscript removes the prime shells outside $S$ from a global kernel on $C_F$ via the shell differences (11.3.1.4), calls the residual global kernel $k_S$, and invokes Theorem~12.4.1.1 to assert
\[
  \int_{C_F} k_S(x,x)\dd\bar\mu(x)
  = D_S(\varphi) - M\varphi(|\cdot|^{1/2}) - M\varphi(|\cdot|^{-1/2}).
\]
Yet Theorem~12.4.1.1 produces $D_S(\varphi)$ alone, without either subtracted Mellin term. For any nonzero, nonnegative radial function supported in $1 < |g| < 2$---which vanishes on $C_F^1$ and satisfies the hypotheses of Theorem~12.5.1.1---both $M\varphi(|\cdot|^{1/2})$ and $M\varphi(|\cdot|^{-1/2})$ are strictly positive, so the discrepancy does not vanish. Equation~(11.3.1.4) merely defines telescopic shell differences; even granting that discarded shells have zero trace leaves $\int_{C_F} k_S(x,x)\dd\bar\mu(x)$ uncomputed. Section~12.3 (p.~514), assigned to compare $k_S$ on $C_F$ with the cutoff kernels on $J_{F,S}$, consists solely of a heading and a footnote to insert the section.

\textbf{Why the Chapter~12 trace formulas have full GRH strength.} The two missing polar terms themselves have a familiar representation-theoretic interpretation: a two-dimensional quotient on which dilations act by $|\cdot|^{1/2}$ and $|\cdot|^{-1/2}$ contributes precisely their sum to the convolution trace (cf.\ \href{https://arxiv.org/pdf/math/0412277}{Meyer, \emph{A spectral interpretation for the zeros of the Riemann zeta function}}, pp.~1--2, for the full zero divisor of $\zeta(s)$). The far deeper obstacle is strip-width invariance: Theorem~12.5.1.1 asserts that the \emph{same} geometric and polar terms equal the zero sum inside \emph{every} centered vertical strip, however narrow.

As proved in Appendix~\ref{app:strip-grh}, this family of strip-invariant trace formulas is already logically equivalent to GRH. Subtracting the identity at a narrower strip width from that at a wider width (for the same test function $\varphi$) cancels the geometric and polar right-hand sides, leaving only the sum over zeros in the intermediate vertical band. Translating $\varphi$ in log-norm turns each surviving zero into an exponential in the translation parameter, so taking a Laplace transform converts the vanishing difference into an identically zero meromorphic function whose poles sit at those zeros (with nonzero residues for a suitably phase-adjusted test). No zero can therefore lie in the band; comparing arbitrarily narrow strips with Weil's classical explicit formula for the full zero sum forces GRH, and conversely GRH implies that every positive-width centered strip already captures the full zero sum. Note that this equivalence applies to the numerical trace formulas assembled in Theorem~12.5.1.1, for which the missing kernel comparison of \S12.3 is a key input.

\subsection{Issue 3: Bounding zero multiplicities requires a separate capped line identity}
\msloc{Expectation 1.7.1; Theorem 12.5.1.1, pp.~519--520; Theorems 13.3.1.1--13.3.1.2, p.~522; \S14.1.2}
\issuefinding{False restatement; unproved and unwritten theorems.}{Presenting Theorem~13.3.1.1 as a restatement of Theorem~12.5.1.1 is false: it alters the test-function hypothesis, restricts the sum to the critical line, and caps each zero multiplicity $m$ at $\min(m,d)$, none of which follows from Chapter~12's strip formula. Both this capped critical-line identity and Theorem~13.3.1.2 (which ends at its title) remain unproved.}

Chapter~12 sums zeros in a positive-width strip with their \emph{full} multiplicities against tests vanishing on $C_F^1$; by contrast, Theorem~13.3.1.1 (p.~522) imposes Mellin moment conditions in place of $\varphi|_{C_F^1}=0$, restricts the sum to the critical line, and caps multiplicities at $\min(m,d)$ where $d=[F:\Q]$. Parameterize idele classes as $g = \tau(e^u)k$ ($k\in C_F^1$), and for a unitary character $\chi$ of $C_F$ with restriction $\omega = \chi|_{C_F^1}$, write $\chi(\tau(e^u)) = e^{ix_0(\chi)u}$ and $f_\omega(u) = \int_{C_F^1}\varphi(\tau(e^u)k)\omega(k)\dd k$. In these coordinates, the integral in (13.3.1.2) is $\int_{\R} e^{(s+ix_0(\chi))u} f_\omega(u)\dd u$; requiring its derivatives in $s$ of orders $0 \le j < d$ to vanish at $s = -ix_0(\chi)$ imposes precisely the moment conditions
\[
  \int_{\R} u^j f_\omega(u)\dd u = 0, \qquad 0 \le j < d.
\]
Already over $\Q$ ($d=1$, trivial $\omega$), the conditions $f(0)=0$ (Chapter~12) and $\int_{\R} f(u)\dd u = 0$ (Chapter~13) are incomparable: a positive bump in $u>0$ satisfies only the first, while a bump positive at $0$ minus a disjoint bump of equal integral satisfies only the second.

Extending a character $\omega$ of $C_F^1$ trivially along $\tau$, a critical-line point $\omega|\cdot|^{1/2+it}$ receives Mellin weight $\int_{\R} f_\omega(u)e^{itu}\dd u$, which vanishes at $t=0$ by the $j=0$ moment condition. (For finite-order characters, $t=0$ corresponds canonically to the central point $s=1/2$, independently of $\tau$.) That central blindness is \emph{consistent} with the manuscript's scope, which expressly exempts a possible central zero from its multiplicity bound (Expectation~1.7.1 and \S14.1.2). Moreover, away from $t=0$ the moment-zero test space \emph{would} suffice to prove $m_t\le d$ \emph{if} one had both GRH and the capped line identity: taking $f = g * \widetilde{g}$ with $g = h^{(d)}$, $h\in C_c^\infty(\R)$, and $\widetilde{g}(u)=\overline{g(-u)}$ (whose Fourier transform $t^{2d}|\widehat{h}(t)|^2$ vanishes to order $2d$ at $t=0$, so $f$ satisfies all $d$ moment conditions) and subtracting the capped identity from Weil's explicit formula yields
\[
  \sum_t \bigl(m_t - \min(m_t, d)\bigr)\, t^{2d}|\widehat{h}(t)|^2 = 0.
\]
Nonnegativity and the freedom to choose $\widehat{h}(t)\ne 0$ for any $t\ne 0$ then force $m_t\le d$.

\textbf{Why this remains a fatal gap for multiplicity.} What is missing is any mechanism producing the capped line identity itself. Every positive-width strip trace in Chapter~12 counts an on-line zero of multiplicity $m_t$ with full weight $m_t$, however narrow the strip; shrinking strips cannot turn $m_t$ into $\min(m_t,d)$ (which for $F=\Q$ and $\chi=1$ would imply the \emph{simplicity of every non-central zero of $\zeta(s)$}, a conjecture well beyond RH). Theorem~13.3.1.2 on p.~522, meant to establish the divisor result, ends immediately after its title with footnotes requesting its statement and proof.

\section{Substantial Unproved Theorems on the Spectral Side}
\label{sec:spectral}

Chapters~9 and~10 build the spectral strip quotient on weighted $L^2$ (``Hardy--Sobolev'') spaces: in $x=\log|g|$, primal spaces carry growing weights $(1+x^2)^{s/2}\dd x$ (times $e^{2\alpha|x|}$ on a strip of half-width $\alpha>0$), while dual spaces carrying Mellin point evaluations have reciprocal decaying weights $w_s(x) = (1+x^2)^{-s/2}$ (times $e^{-2\alpha|x|}$). Two errors here leave major functional-analytic theorems unproved; even resolving both would leave intact the fatal global obstructions of Section~\ref{sec:fatal}.

\subsection{Issue 4: Reversing the decreasing projections leaves an open strip approximation theorem}
\msloc{Theorem 9.3.2.6, p.~426; applied in Theorem 10.2.2.1, p.~446 (\S10.2.2)}
\issuefinding{False theorem; unproved consequence.}{Theorem~9.3.2.6 is false: its proof reverses the inclusion of the height-$T$ zero-kernel subspaces ($V_{T'}\subseteq V_T$ for $T'>T$, not $V_T\subseteq V_{T'}$), so the orthogonal projections $P_T$ onto them fail even to converge weakly to the identity. This leaves Theorem~10.2.2.1 unproved, resting on an unproved strip-norm approximation theorem for the zero kernel.}

In the manuscript's Hilbert space of holomorphic Mellin transforms on an open strip, point and derivative evaluations at interior points are continuous. Given a fixed function $\phi$ whose Mellin transform prescribes a zero divisor, Theorem~9.3.2.6 defines $V_T$ as the closed subspace of functions vanishing (with required derivatives) at all zeros of height at most $T$, and lets $P_T$ be the orthogonal projection onto $V_T$. The proof claims $V_T \subseteq V_{T'}$ for $T < T'$ and concludes that $P_T\to I$ strongly as $T\to\infty$. In fact, raising $T$ adds linear constraints, so the inclusion runs the opposite way: $V_{T'}\subseteq V_T$. On $G=\R_{>0}$ and a strip containing $0$, choosing $\phi$ with Mellin transform $z e^{z^2}$ (a differentiated Gaussian in $x$, with a single simple zero at $0$) makes $V_T = \{F : F(0)=0\}$ for all $T>0$. Every $P_T$ then annihilates the value at $0$ of the undifferentiated Gaussian's transform $e^{z^2}$, ruling out even weak convergence to $e^{z^2}$.

\textbf{Strategic burden of a repair.} Write $A$ for the closed geometric subspace (generated by adelic theta/zeta transforms) divided out to form the spectral quotient, and $D = \bigcap_{T>0} V_T$ for the common kernel of all prescribed zero and derivative evaluations. Passing to the complementary projections $I - P_T$ recovers the closed span of the evaluation representers, namely $D^\perp$, whereas the dual of the intended quotient is $A^\perp$. Equating them in Theorem~10.2.2.1 requires proving $A = D$ in the completed strip norm---a nontrivial spectral-synthesis or approximation theorem asserting that every strip function vanishing on the zero divisor is approximable by the specified geometric test functions. Correcting the direction of $P_T$ supplies no proof of $A = D$.

\subsection{Issue 5: Clustered simple \texorpdfstring{$\zeta$}{zeta}-zeros defeat coordinatewise trace-norm bounds}
\msloc{Proposition 10.2.1.3, p.~443; individual-coordinate bound after (10.3.1.18), p.~451; Theorem 10.3.2.1, p.~453}
\issuefinding{False bound; unproved theorem.}{The uniform bound on individual zero-evaluation/biorthogonal pairs asserted on p.~451 (after (10.3.1.18)) is false for actual simple critical-line zeros of $\zeta(s)$, whose clustering forces the biorthogonal functional norms to blow up. This invalidates the proof of Theorem~10.3.2.1, leaving spectral nuclearity on the strip unproved.}

To prove spectral nuclearity (Theorem~10.3.2.1), the manuscript expands the operator in zero-evaluation vectors and biorthogonal functionals, relying on a uniform bound on individual biorthogonal pairs (p.~451). Take $F=\Q$, the trivial character, the critical line $\RePt z = 1/2$, and a fixed Sobolev index $s>1$; a simple zero $\frac{1}{2}+i\gamma$ gives the evaluation vector $\eta_\gamma(x) = e^{i\gamma x}$ with constant dual norm $\|\eta_\gamma\|^2 = \int_{\R} w_s(x)\dd x$. By the Riemann--von Mangoldt formula and \href{https://arxiv.org/pdf/1002.4127}{Bui--Conrey--Young, Theorem~1.1}, $(0,T)$ contains $\gg T\log T$ distinct simple critical-line ordinates, whereas any uniformly separated set of points in $(0,T)$ has size at most $O(T)$. Hence there exist pairs of distinct simple critical-line ordinates with $|\gamma-\gamma'|\to 0$, for which dominated convergence gives
\[
  \|\eta_\gamma - \eta_{\gamma'}\|^2 = \int_{\R} |e^{i(\gamma-\gamma')x} - 1|^2 w_s(x)\dd x \longrightarrow 0.
\]
Since the biorthogonal functional $b_\gamma$ takes values $1$ at $\eta_\gamma$ and $0$ at $\eta_{\gamma'}$, we have $1 \le \|b_\gamma\|\,\|\eta_\gamma - \eta_{\gamma'}\|$, forcing $\|b_\gamma\|\to\infty$. Rescaling preserves the product $\|b_\gamma\|\|\eta_\gamma\|$, and inserting the extra decaying weight $e^{-2\alpha|x|}$ for a centered strip of half-width $\alpha>0$ leaves the argument unchanged.

\textbf{Prospect of repair.} Unbounded individual coordinate functionals do not by themselves rule out a finite trace norm for the \emph{sum}: in a $2\times 2$ model with eigenvectors $e_1$ and $e_1 + \delta e_2$ and eigenvalues $m(t), m(t+\delta)$ for a smooth function $m$, the individual spectral projections grow like $\delta^{-1}$, yet in orthonormal coordinates $(e_1,e_2)$ the off-diagonal entry is the bounded divided difference $(m(t+\delta)-m(t))/\delta$. On the critical line, \href{https://arxiv.org/pdf/math/9811068}{Connes (\S III, Cor.~2 and App.~I)} proved a trace-class convolution theorem (with Sobolev-truncated multiplicities) by grouping frequencies into intervals. Because a nonzero holomorphic function on an open strip cannot vanish on a vertical interval inside the strip, Connes's frequency cutoff has no literal strip analogue; a genuine repair would need to group, bound, and sum all zero clusters directly in the strip norm.

\section{Locally Reparable Errors and Bypassed Deductions}
\label{sec:reparable}

The remaining five issues falsify specific theorems or deductions as printed, but each either admits a plausible local correction or would be bypassed if the core obstructions of Section~\ref{sec:fatal} were resolved.

\subsection{Issue 6: Failed overlap-gluing step in the Global Converse Theorem}
\msloc{Theorem 6.3.2.1 and eq.~(6.3.2.11), pp.~304--306; applied in \S10.2.1}
\issuefinding{False gluing step; unproved theorem.}{In the proof of Theorem~6.3.2.1, the claim after (6.3.2.11) that the regularized right- and left-half-plane pieces agree on $0<\RePt s<1$ is false for the trivial character: their difference is $(1-2s)$ times the global zeta integral $Z(f,s)$, which is not identically zero even when $f(0)=\widehat{f}(0)=0$. This leaves Theorem~6.3.2.1's image and closed-range characterization unproved for $\zeta(s)$.}

To identify the global zeta-transform image used in \S10.2.1, Theorem~6.3.2.1 constructs holomorphic functions on right and left half-planes from $f$ and $\widehat{f}$ and applies a difference map $R$ on $0<\RePt s<1$ to glue them via $\ker R$. For $F=\Q$, the trivial character, and finite factors $\mathbf{1}_{\Z_p}$, multiplying by the missing Euler factors produces the global zeta integrals $Z(f,s)$ and $Z(\widehat{f},1-s)$, which agree on $0<\RePt s<1$ by Tate's functional equation. To clear poles, however, (6.3.2.11) multiplies the first piece by $1-s$ and the second by $s$, so their difference on the overlap is
\[
  (1-s)Z(f,s) - s Z(\widehat{f},1-s) = (1-2s)Z(f,s).
\]
Even on the zero-moment subspace $f(0)=\widehat{f}(0)=0$, the choice $f_\infty(x) = e^{-\pi x^2} - 3e^{-4\pi x^2} + 2e^{-16\pi x^2}$ satisfies $f_\infty(0) = 1-3+2=0$ and $\widehat{f}_\infty(0) = 1 - \tfrac{3}{2} + \tfrac{2}{4} = 0$, while $Z(f,s)$ is a nonzero multiple of
\[
  \pi^{-s/2}\Gamma(s/2)\zeta(s)\bigl(1 - 3\cdot 2^{-s} + 2\cdot 4^{-s}\bigr),
\]
so $(1-2s)Z(f,s)$ does not vanish on $0 < \RePt s < 1$.

\textbf{Prospect of repair.} With $F(s)=(1-s)Z(f,s)$ and $G(s)=sZ(\widehat{f},1-s)$, one can fix the algebraic gluing error by replacing naive equality on the overlap with $sF(s) = (1-s)G(s)$---or by dividing out the linear factors to recover a single meromorphic function whose poles are removable on the zero-moment space. The asserted image characterization and closed-range property in the specified topology still require proof.

\subsection{Issue 7: Inverted Fourier cutoff level in a Paley--Wiener orbital quotient}
\msloc{Proposition 7.3.3.2 and Definition 7.3.3.3, pp.~358--359; Theorem 7.3.3.4, p.~360}
\issuefinding{False containment (ill-defined quotient).}{Definition~7.3.3.3 requires $\mathrm{PV}_p(\mathrm{PW}_{X,0})\subseteq\mathrm{PW}_{X/p,0}$ to form $\mathrm{PW}_{X/p,0}/\mathrm{PV}_p(\mathrm{PW}_{X,0})$, which is false: $\mathrm{PV}_p$ genuinely enlarges Fourier support to $[-pX,pX]$. The quotient in Definition~7.3.3.3 is therefore undefined, invalidating Theorem~7.3.3.4's basis construction as stated.}

For $F=\Q$, $S=\{\infty\}$, one added prime $p$, and the trivial character, $\mathrm{PW}_{X,0}$ consists of $\{\infty,p\}$-zeta classes admitting a smooth real representative $f$ with $\supp\widehat{f}\subset[-X,X]$ and $f(0)=\widehat{f}(0)=0$. Proposition~7.3.3.2 represents $\mathrm{PV}_p$ by $f(x)\mapsto f(x)-f(px)$, mapping level $X$ into level $pX$, yet Definition~7.3.3.3 forms the quotient $\mathrm{PW}_{X/p,0}/\mathrm{PV}_p(\mathrm{PW}_{X,0})$ (on which Theorem~7.3.3.4 builds a basis). Choose a nonzero smooth even function $h(\xi)$ supported in $X/p<|\xi|<X$ with $\int h=0$, and set $f=\mathcal{F}^{-1}h$; then $f(0)=\widehat{f}(0)=0$, so $f$ represents a class in $\mathrm{PW}_{X,0}$. Its image $f(x)-f(px)$ has Fourier transform
\[
  h(\xi) - p^{-1}h(\xi/p).
\]
Whenever $h(\xi_0)\ne 0$, the point $\xi=p\xi_0$ satisfies $|\xi|>X$, so $h(p\xi_0)=0$ and the image transform evaluates to $-p^{-1}h(\xi_0)\ne 0$ outside $[-X,X]$. Any alternative representative can be replaced by its even part without enlarging its Fourier support, and two even representatives of the same class share real Mellin transforms after cancelling $(1-p^{-s})^{-1}$. By Mellin uniqueness on $\R_{>0}$, no representative can lie at level $X/p$, so the quotient in Definition~7.3.3.3 is not well defined.

\textbf{Prospect of repair.} Reindexing the numerator to $\mathrm{PW}_{pX,0}$ fixes the failed containment and still exhausts the inductive family for fixed $p$, provided one tracks the basis and cutoff shifts as primes are added. However, whenever Chapter~11's log-norm truncation and closure are subsequently performed, the fatal shell collapse of Issue~1 persists.

\subsection{Issue 8: Ambient function support vs.\ distributional restriction on hyperplanes}
\msloc{Theorem 7.3.1.6, pp.~344--345; Theorem 8.3.0.1, pp.~387--389; applied in \S8.4}
\issuefinding{False theorem.}{Theorem~8.3.0.1 is false as stated: asserting that every orbital functional $T$ attains $T(f)=T(f_0)$ on an ambient Schwartz--Bruhat function $f_0$ supported on the coordinate hyperplanes $\bigcup_{v\in S}\{x_v=0\}$ forces $f_0\equiv 0$ and hence $T\equiv 0$, contradicting nonzero local zeta residues from Theorem~7.3.1.6.}

For finite $S$ and a functional $T$ in the orbital dual (including local zeta principal parts by Theorem~7.3.1.6), the open subset $\prod_{v\in S} F_v^\times$ is dense in $\prod_{v\in S} F_v$. Consequently, any continuous ambient function $f_0$ on $\prod_{v\in S} F_v$ supported on $\bigcup_{v\in S}\{x_v=0\}$ vanishes identically, forcing $T(f)=T(f_0)=0$ for all $f$. Yet for $F=\Q$, $S=\{\infty,p\}$, and $f(x_\infty,x_p) = e^{-\pi x_\infty^2}\mathbf{1}_{\Z_p}(x_p)$ (with $\mathrm{vol}(\Z_p^\times)=1$), $Z_S(f,s)$ is proportional to
\[
  \frac{\pi^{-s/2}\Gamma(s/2)}{1 - p^{-s}},
\]
whose residue $T(f) = \Res_{s=2\pi i/\log p} Z_S(f,s)$ is nonzero.

\textbf{Prospect of repair.} While a continuous ambient test \emph{supported on} $\{x_p=0\}$ is zero, restricting an ambient test to $\{x_p=0\}$ does recover individual local residues: if $h\in C_c^\infty(\Q_p)$ equals $h(0)$ on $p^N\Z_p$, that ball contributes $h(0)p^{-Ns}/(1-p^{-s})$ while outer shells give a finite sum of entire functions, yielding residue $h(0)/\log p$ at $2\pi ik/\log p$---an evaluation distribution supported at $0$. Distributional restriction to coordinate hyperplanes can therefore represent single-place residues, though coincident poles across multiple places and comparison with a global diagonal trace still require separate arguments.

\subsection{Issue 9: Identical vanishing of the idelic kernel-lift weight vs.\ partition lifts}
\msloc{Theorem 11.3.2.1 and (11.3.2.1)--(11.3.2.2), p.~495; Proposition 12.2.2.2; applied on p.~520}
\issuefinding{False identity; unproved corollary.}{Identity~(11.3.2.1) is false for the averaging weight $u_0(x,y)$ defined in (11.3.2.2): $u_0$ vanishes identically because its indicator selects a single measure-zero point, whereas the right-hand side is strictly positive for suitable radial bumps over $\Q$. The idelic kernel lift and the large-prime shell vanishing invoked on p.~520 therefore remain unproved.}

Formula~(11.3.2.2) lifts kernels from $C_F$ to $\A_F^\times$ via an averaging weight $u_0(x,y)$ defined by summing over $\xi\in F^\times$ the integral of $\mathbf{1}_{\tau(\R_{>0})}(\xi g x)\mathbf{1}_{\tau(\R_{>0})}(\xi g y)$ over $g$ in a fundamental domain $E_F^1$ of the norm-one ideles. Since $|\xi g x|=|x|$ and $\tau(\R_{>0})$ contains a unique element of norm $|x|$, the first indicator holds only at the single point $g = \xi^{-1}\tau(|x|)x^{-1}$, which has Haar measure zero in the nonatomic norm-one ideles (containing $\prod_p\Z_p^\times$ over $\Q$). Hence $u_0(x,y)=0$ everywhere, contradicting (11.3.2.1) for $F=\Q$, $S=\{\infty\}$, and a nonnegative radial bump positive at norm $p>1$, where the right-hand side is strictly positive.

\textbf{Prospect of repair.} Using the $\Q^\times$-fundamental domain $D = \R_{>0}\times\Zhat^\times$, whose rational translates partition $\A_\Q^\times$, one can construct normalized lifts satisfying $\sum_{r\in\Q^\times} b(rx)=1$, restoring the total trace. Individual shell traces, however, depend on the choice of lift: both $b_0=\mathbf{1}_D$ and $b_1=\tfrac{1}{2}(\mathbf{1}_D + \mathbf{1}_{pD})$ satisfy $\sum_{r\in\Q^\times} b_i(rx)=1$, yet Proposition~12.2.2.2's skew-diagonal weight $\sum_{r\in\Q^\times} b_i(rx)b_i(r\xi x)$ ($x\in D$) puts weight $1$ at $\xi=1$ for $b_0$, versus $\tfrac{1}{2}$ at $1$ and $\tfrac{1}{4}$ at each of $p,p^{-1}$ for $b_1$. Thus the required arithmetic shell allocation---or Chapter~12's narrower large-prime vanishing property---must still be proved for the actual kernel.

\subsection{Issue 10: Divergent archimedean cutoffs in \texorpdfstring{\S1.4}{Section 1.4}, bypassed if strip-invariance holds}
\msloc{Sections 1.4.2--1.4.5, pp.~28--31; bound asserted on p.~29 and repeated on p.~54}
\issuefinding{False bound; invalid deduction.}{The local uniform bound on the archimedean terms asserted in \S1.4.3 (p.~29, repeated on p.~54) is false: for $F=\Q$, the trivial character, and real $z>\max(1,\sup I)$, the archimedean integrals diverge to $+\infty$, invalidating the Montel-limit deduction of GRH in \S1.4.5.}

Sections~1.4.2--1.4.5 test the trace formula against smooth nonnegative cutoffs $\varphi_n(t)$ vanishing near $t=1$ and approaching $t^{-z}\mathbf{1}_{t>1}$ on $t=|g|>0$, relying on this local uniform bound to extract a holomorphic Montel limit (\S1.4.5). Take $F=\Q$, the trivial character, and real $z > \max(1,\sup I)$. Combining Haar measure $dt/t$ with Weil's archimedean kernel $t^{1/2}/|1-t|$ on the positive real branch $1<t<2$, Fatou's lemma gives
\[
  \liminf_{n\to\infty}\int_1^2 \frac{\varphi_n(t)\,t^{-1/2}}{t-1}\dd t
  \;\ge\; \int_1^2 \frac{t^{-z-1/2}}{t-1}\dd t \;=\; +\infty.
\]
Neither principal-value regularization (which vanishes because $\varphi_n\equiv 0$ near $1$) nor the negative real branch (where the integrand is likewise nonnegative) can cancel this divergence.

\textbf{Strategic assessment.} This refutes the printed Montel deduction. Subtracting tests at two $z$-values could cancel the $t=1^+$ singularity, at the cost of controlling the resulting differences in all remaining terms. More fundamentally, a trace identity on a single strip can at most confine zeros to that strip; conversely, if Chapter~12's family of strip-invariant identities were ever established, the translated-bump argument of Issue~2 would already force GRH without invoking \S1.4.

\section{Concluding Remarks}

While Issues~6--10 admit plausible local fixes or workarounds, the central obstructions to the manuscript's program lie in Issues~1--3, reinforced by the open spectral theorems of Issues~4--5. Truncating before taking closures in Chapter~11 collapses every prime shell to zero for $\zeta(s)$, erasing the geometric mechanism offered to keep the trace invariant as the spectral strip narrows; meanwhile, Section~12.3 and Theorem~13.3.1.2 remain unwritten precisely where the GRH-strength strip comparison and the capped critical-line multiplicity identity would have to be proved.

\clearpage
\appendix
\section{Why the Chapter 12 strip formulas imply GRH}
\label{app:strip-grh}

After rearranging Theorem~12.5.1.1, one side sums over the nontrivial zeros $\rho$ inside the centered strip $|\RePt\rho - \tfrac12| < a$ (for any half-width $a>0$), while the other side equals the sum of the two polar terms minus the local distributions---an expression independent of $a$ for any fixed test function. We show below that whenever this identity holds at two half-widths, every zero in the wider strip must also lie in the narrower one; comparing arbitrarily narrow strips with Weil's explicit formula for the full zero sum then yields GRH.

\subsection*{Selecting one Hecke character}

Fix a number field $F$, and parameterize idele classes as $g = \tau(e^u)k$ with $u = \log|g|\in\R$ and $k$ in the compact norm-one subgroup $K = C_F^1$, equipped with Haar measure $du\dd k$ normalized by $\int_K \dd k = 1$. Fix a unitary character $\omega$ of $K$ and extend it trivially along $\tau$ via $\omega(\tau(e^u)k) = \omega(k)$. Varying $s\in\C$ gives the family of quasi-characters $c = \omega|\cdot|^s$ restricting to $\omega$ on $K$, with associated Hecke $L$-function $s \mapsto L(\omega|\cdot|^s)$. Write $\rho$ for any nontrivial zero of this $L$-function, $w = \rho - \tfrac12$ for its centered coordinate, and $m_w$ for its multiplicity. In these coordinates GRH asserts that $\RePt w = 0$ for every such zero; choosing a different unitary extension of $\omega$ merely shifts $s$ by a purely imaginary constant and leaves the statement unchanged.

Choose a smooth complex-valued function $h\in C_c^\infty((0,1))$, and for each real $x\ge 0$ define the test function
\[
  \varphi_x(\tau(e^u)k) = \overline{\omega(k)}\,h(u-x).
\]
Its log-norm support lies in $(x, x+1)$, so $\varphi_x$ vanishes on $K$ as Chapter~12 requires. For each fixed $x\ge 0$, enlarging the finite set $S$ in the proof of Theorem~12.5.1.1 to include all finite places of norm at most $e^{x+1}$, together with the places required for $\omega$ and $F$, produces a single set $S$ valid for both strip widths at that $x$.

To see which zeros $\varphi_x$ detects, let $\eta$ be any character of $K$, extended trivially along $\tau$. By character orthogonality, the angular integral $\int_K \overline{\omega(k)}\eta(k)\dd k$ equals $1$ for $\eta = \omega$ and $0$ otherwise, so only the $\omega$-component contributes to the zero sum. At a zero $c = \omega|\cdot|^\rho$ in this component, Chapter~12 evaluates the Mellin transform at $c|\cdot|^{-1/2} = \omega|\cdot|^w$, assigning it the weight
\[
  M\varphi_x(\omega|\cdot|^w)
   = \int_{\R} h(u-x)e^{wu}\dd u
   = e^{wx}H(w), \qquad H(w) = \int_0^1 h(v)e^{wv}\dd v,
\]
after substituting $v = u-x$. Translating a single test function in log-norm thus varies the exponential factor $e^{wx}$ while keeping $H(w)$ fixed.

\subsection*{Subtracting the two strip formulas}

Applying Theorem~12.5.1.1 at half-width $a$ to $\varphi_x$ yields
\[
  \sum_{|\RePt w|<a} m_w H(w)e^{wx}
   = M\varphi_x(|\cdot|^{1/2}) + M\varphi_x(|\cdot|^{-1/2}) - D_S(\varphi_x),
\]
where each distinct centered zero $w$ in the $\omega$-component appears once with its multiplicity $m_w$. Fix $0 < \alpha < \beta$ and subtract the identity at half-width $\alpha$ from that at half-width $\beta$. For each fixed $x\ge 0$, the right-hand sides cancel, leaving on the left only the set of surviving zeros $B = \{w : \alpha \le |\RePt w| < \beta\}$ (which includes the inner boundary $|\RePt w| = \alpha$ because the narrower strip is open):
\begin{equation}
  \sum_{w\in B} m_w H(w)e^{wx} = 0 \qquad (x \ge 0).
  \label{eq:band-translates}
\end{equation}
This series is absolutely convergent: Weil's classical explicit formula (Theorem~B.4.4.7 of the manuscript) guarantees absolute convergence of the full zero sum for any smooth compactly supported test, which at $x=0$ gives $\sum_{w\in B} m_w|H(w)| < \infty$, and for any $x\ge 0$ translation increases the absolute sum by at most a factor of $e^{\beta x}$.

\needspace{14\baselineskip}
\subsection*{Detecting one zero}

For a finite sum one could now appeal to the linear independence of distinct exponentials; to handle the infinite sum, take the Laplace transform in $x$. For any $z\in\C$ with $\sigma = \RePt z > \beta$, multiply \eqref{eq:band-translates} by $e^{-zx}$ and integrate over $x\ge 0$. Termwise integration is justified by Fubini's theorem via the bound
\[
  \sum_{w\in B} m_w|H(w)|\int_0^\infty e^{-(\sigma-\RePt w)x}\dd x
  \;\le\; \frac{1}{\sigma-\beta}\sum_{w\in B} m_w|H(w)| < \infty.
\]
Since $\int_0^\infty e^{-(z-w)x}\dd x = 1/(z-w)$, we obtain
\[
  G_h(z) := \sum_{w\in B}\frac{m_w H(w)}{z-w} = 0 \qquad (\RePt z > \beta).
\]
Because $B$ is a discrete subset of $\C$ (being a set of zeros of a fixed Hecke $L$-function), every compact subset of $\C\setminus B$ remains at a positive distance from $B$. The absolute summability $\sum_{w\in B} m_w|H(w)| < \infty$ therefore yields uniform convergence on compact subsets of $\C\setminus B$, so $G_h$ is holomorphic on $\C\setminus B$. Because this punctured domain is connected (any path can detour around the finitely many points of $B$ in a bounded region), the identity theorem extends $G_h \equiv 0$ from the half-plane $\RePt z > \beta$ to all of $\C\setminus B$.

Suppose now that $w_0\in B$, and choose a small disk about $w_0$ containing no other element of $B$. Removing the $w_0$ term leaves a series that converges uniformly on a smaller closed disk by the same distance bound and is therefore holomorphic at $w_0$, giving
\[
  0 = \lim_{z\to w_0}(z-w_0)G_h(z) = m_{w_0}H(w_0).
\]
Since this holds for every $h\in C_c^\infty((0,1))$, choose a smooth nonnegative bump $b\in C_c^\infty((0,1))$ that is not identically zero and set $h(v) = e^{-i(\ImPt w_0)v}b(v)$. This phase factor cancels the oscillation in the Mellin integral at $w_0$ and yields
\[
  H(w_0) = \int_0^1 b(v)e^{(\RePt w_0)v}\dd v > 0.
\]
Because $m_{w_0}\ge 1$, this contradicts $m_{w_0}H(w_0)=0$; hence $B$ is empty.\footnote{For the Riemann zeta function, one may even force both polar terms to vanish separately by replacing the phase-adjusted test $h$ above with $h''-\tfrac14 h$. Two integrations by parts multiply its Mellin transform by $w^2-\tfrac14$, which vanishes at the polar points $w=\pm\tfrac12$ yet remains nonzero at every centered nontrivial zero (since $|\RePt w|<\tfrac12$).}

\subsection*{Passing from one band to GRH}

Weil's classical explicit formula provides the outer comparison unconditionally: since every nontrivial zero satisfies $0 < \RePt\rho < 1$, the full zero sum is precisely the sum for $\beta = \tfrac12$, and its geometric and polar terms agree with Chapter~12 in the stated Haar normalization. If an off-line zero $w_0$ existed (so $0 < |\RePt w_0| < \tfrac12$), choosing any asserted half-width $\alpha \in (0, |\RePt w_0|)$ would place $w_0$ in the band $\alpha \le |\RePt w| < \tfrac12$. Applying the preceding argument to Theorem~12.5.1.1 at half-width $\alpha$ and Weil's explicit formula at $\beta = \tfrac12$ shows that this band is empty---a contradiction. Since $F$ and $\omega$ were arbitrary, GRH follows for every Hecke $L$-function.

Conversely, under GRH every positive-width centered strip already contains all nontrivial zeros with their original multiplicities, so Weil's explicit formula immediately yields the numerical identity of Theorem~12.5.1.1 at every half-width $a>0$. This equivalence governs only the locations of the zeros; bounding their multiplicities would require the separate capped critical-line identity discussed in Issue~3.

\end{document}
